\documentclass[11pt]{article}
\usepackage[margin=1in]{geometry}
\usepackage{times}
\usepackage{url}
\pagestyle{empty}
\begin{document}

\vspace{-2ex}
\noindent\textbf{Cover Letter}\\[1ex]
\noindent\textbf{To:} Editor-in-Chief, ACM Transactions on Computing Education\\
\textbf{From:} Wenbin Hu, School of Computer Science, Xijing University\\
\textbf{Date:} July 30, 2026\\
\textbf{Subject:} Manuscript Submission --- ``Automated Programming Verdict Classification from Submission Metadata: Comparing Machine Learning and Large Language Models in Competitive Programming''\\[2ex]

\noindent Dear Editor,\\[1ex]

We are pleased to submit our manuscript, ``Automated Programming Verdict Classification from Submission Metadata: Comparing Machine Learning and Large Language Models in Competitive Programming,'' for consideration in ACM Transactions on Computing Education.\\[1ex]

\noindent\textbf{Why This Paper Fits TOCE.} Automated feedback is central to computing education, yet most existing tools require source code access. We address a deceptively simple question: \textbf{what verdict information can be recovered from metadata alone?} Our answer reveals a platform-encoding boundary --- a natural limit that neither ML nor LLMs can overcome without source code. The paper provides educators with cost-effective verdict-screening tools (2\,ms inference, \$0.01 per 1,000 classifications) and an empirical framework for comparing supervised ML and zero-shot LLM approaches.\\[1ex]

\noindent\textbf{Novelty and Contributions.}
\begin{itemize}\setlength\itemsep{0.4em}
\item \textbf{First controlled ML--LLM comparison} for programming verdict classification from metadata alone, with identical protocols across three ML models (Random Forest, Gradient Boosting, Logistic Regression) and two LLM configurations (DeepSeek-V3, Qwen2.5:3B).
\item \textbf{The platform-encoding boundary}: CE, TLE, and MLE (18.4\% of submissions) are deterministically recoverable from execution metadata (F1 $\ge$ 0.89) because platform measurement thresholds define these outcomes. The WA$\leftrightarrow$RE boundary is not metadata-separable (RE F1 = 0.52), indicating where source-code analysis may be required.
\item \textbf{Statistical and deployment rigor}: Holm-Bonferroni-corrected McNemar tests, ablation analysis, Gini importance, leave-one-user-out cross-validation (82.5\% accuracy), and cross-difficulty generalization (71--94\% across problem difficulties).
\item \textbf{Open dataset and code}: \url{https://github.com/huwenbin123huwenbin/programming-error-classification}
\end{itemize}

\noindent\textbf{Key Findings.} Gradient Boosting (95.02\%) outperforms DeepSeek-V3 zero-shot (76.65\%) by 18.4 pp; even with identical features GB (92.0\%) surpasses DeepSeek-V3 by 15.4 pp. Execution time is the dominant predictor (8.8 pp accuracy drop when removed). We note that this comparison is asymmetric: ML models were trained on 10,688 labeled samples while LLMs received no task-specific training.\\[1ex]

\noindent\textbf{Relevance to TOCE Scope.} This work speaks to TOCE's interest in: (1) automated assessment tools for resource-constrained settings; (2) feedback automation characterizing which verdicts benefit from immediate feedback (CE, TLE, MLE) and which require human expert review; (3) empirical computing education research with transparent statistical analysis.\\[1ex]

\noindent\textbf{Submission Materials.} Manuscript: 28 pages, 8 tables, 7 figures. Highlights: 5 key contributions. Supplementary: per-class F1 scores, Cohen's $h$ effect sizes, balanced accuracy, confusion matrices, ROC/PR curves.\\[1ex]

We confirm this manuscript is original, has not been published elsewhere, and is not under review at any other venue.\\[2ex]

Best regards,\\[2ex]

Wenbin Hu\\
School of Computer Science, Xijing University\\
wenbin.hu2026@outlook.com
\end{document}
